% ch3-design.tex -- part of thesis_v3.tex; do not add \documentclass here.
% Draft with the thesis-chapter workflow. Unsupported claims must use \fillin{...}.
% Implementation traceability (file:line) is held in the Implementation Index appendix,
% not in the body of this chapter.

\chapter{Design and Implementation}
\label{ch:design_implementation}

This chapter describes the framework as it was built and configured. It covers the construction and
normalization of the multi-frequency dataset, the multi-modal variational autoencoder that serves as
the power-integrity surrogate, the latent-space optimization stage that searches for decoupling
capacitor placements, the experimental setup, and the active-learning loop that extends the training
set with newly simulated layouts. The subject is the method and its implementation; results are
reported in the following chapter and are not anticipated here.

Notation is fixed as follows. A placement is a binary vector $\mathbf{b} \in \{0,1\}^{52}$ over the
52 decap slots of the board, the decap budget is $K = \sum_i b_i$, and the power-integrity response
is the input impedance magnitude $Z(f)$ sampled on 231 frequency bins spanning 1--600~MHz. The
target impedance mask against which a placement is judged is $Z_{\mathrm{target}}(f)$. The spatial
power-integrity distribution at a single inspection frequency is a heatmap $H$ on a
$64 \times 64$ grid. \fillin{describe the board under study (H-shape test vehicle) in one paragraph
here: outline, stackup, plane pair and port definition. The active-learning configuration names the
ECADSTAR design file H-shape.erf with powerbus Power\_GND and port IC1\_Port1, but no committed
artifact documents the geometry itself.}

\section{Data Preparation and Reprocessing}
\label{sec:data_prep_reprocessing}

\subsection{Raw Power-Integrity Simulation Outputs and File Structure}
\label{subsec:raw_pi_outputs}

Ground truth is produced by ECADSTAR power-integrity simulation. A processing stage converts the raw
export into a layout store made of three parallel trees --- per-design layout arrays, per-sample
heatmaps and per-sample spectra --- together with a manifest index (\texttt{manifest.csv}).

Per design the store holds an occupancy vector of shape $(52,)$ and an impedance spectrum of shape
$(231,)$. These two shapes are the fixed interface of the whole framework: the occupancy graph
encoder is built over 52 slots and the impedance graph encoder over 231 bins. A manifest row is one
(design, inspection frequency) pair, so a layout appears once per anchor frequency, each time with
its own spatial heatmap, while occupancy and the full spectrum are properties of the layout alone.

\fillin{describe the ECADSTAR simulation setup behind the raw export: solver settings, excitation,
meshing, and how the spatial heatmap is exported. Only the active-learning batch invocation of the
simulator is documented; the original bulk data generation is not described in any committed
artifact.}

\fillin{confirm the stored heatmap resolution as a dataset property. The active-learning
configuration sets a heatmap grid size of 64 and the decoder's final conditioning stage operates at
$64 \times 64$, but the dataset metadata records no heatmap shape field, so the resolution is
inferred from the model and the pipeline configuration rather than read from the dataset.}

\figplace{Data flow from the ECADSTAR export through the processing stage to the layout store (layout arrays, heatmaps, spectra, manifest).}

\subsection{Multi-frequency Dataset Construction and Normalization}
\label{subsec:multifreq_dataset}

The normalized training set used throughout this work is
\texttt{datasets/data\_multifreq\_train\_norm\_unbounded}. Its metadata records the stage as
normalized, names \texttt{pipelines/normalize/multifreq.py} as the generating script, and carries the
generation timestamp 2026-07-06T13:39:03.768907+00:00. The set contains 582997 manifest rows over
24979 unique layouts and 24979 layout directories, with 582997 heatmap files and 582997 spectrum
files.

The dataset is multi-frequency in the sense that each layout is simulated at a set of anchor
inspection frequencies. Twenty-four anchors are recorded, in MHz: 10.0, 63.0, 80.0, 100.0, 120.0,
150.0, 170.0, 180.0, 200.0, 230.0, 250.0, 270.0, 280.0, 300.0, 330.0, 350.0, 370.0, 390.0, 400.0,
420.0, 430.0, 450.0, 470.0 and 500.0. Coverage is not uniform. Every anchor carries 24979 samples
except 80~MHz, which carries 8480.
\fillin{state whether the reduced 80~MHz coverage (8480 versus 24979) is deliberate or the residue of
an interrupted simulation batch.}
The anchor set is reported inconsistently across the project's internal documentation, which names
sets of sixteen and of six anchors; the twenty-four values above are those recorded in the metadata
of the dataset actually consumed at training time, and are the ones used here.
\fillin{reconcile the anchor sets recorded in dataset\_meta.json (24), docs/dataset.md (16) and
docs/normalization-and-losses.md (six), and correct the two documents, which predate exp059.}

Before statistics are fitted, the build applies a decap-budget filter at $K \le 30$. The budget is
recovered per design as the number of occupancy entries above 0.5, and layouts exceeding the limit
are dropped both from the statistics fit and from the written output.

Three modalities are normalized separately.

\paragraph{Heatmaps.} The normalization mode is robust, per-MHz, applied in $\log(1+x)$ space and
unbounded. The forward and inverse transforms are
\begin{equation}
z \;=\; \frac{\operatorname{log1p}(\Omega) - \mathrm{median}_{\mathrm{MHz}}}{\mathrm{IQR}_{\mathrm{MHz}}},
\qquad
\Omega \;=\; \mathrm{expm1}\!\left(z \cdot \mathrm{IQR}_{\mathrm{MHz}} + \mathrm{median}_{\mathrm{MHz}}\right),
\label{eq:heatmap_norm}
\end{equation}
where $\Omega$ is the physical impedance of a heatmap pixel in ohms and the median and interquartile
range are fitted per anchor frequency. Two entries illustrate the spread: at 10.0~MHz the median is
0.05887173116207123 and the IQR is 0.03866304084658623, while at 500.0~MHz the median is
1.3578120470046997 and the IQR is 0.6529430150985718. The per-anchor statistics therefore differ by
more than an order of magnitude across the anchor set. Global quality-control references are stored
alongside them: a lower clip of $-1.6791870594024658$, an upper clip of $6.275787830352783$, a
maximum $z$ of $11.23223876953125$ and a background value of $-3.1792$. In the unbounded mode these
clip values serve as references and not as training bounds. The experiment configuration mirrors the
same constants and does not clip reconstructions to that band.

Foreground and background are separated by a physical threshold of 0.05~$\Omega$. Robust clip
percentiles of 0.5 and 99.5 are defined in the same normalization stage but take effect only in the
bounded mode. Two build-time switches select between these behaviours: robust per-MHz statistics are
enabled by default, the unbounded transform is not, so the build that produced this dataset selected
the unbounded mode explicitly.

\paragraph{Impedance spectra.} The 231-bin spectrum is normalized by a single global log z-score with
mean $-1.039238452911377$ and standard deviation $1.7831873893737793$, fitted over 24979 layouts. The
resulting range spans $-2.7170498371124268$ to $2.474900722503662$.

\paragraph{Inspection frequency.} The scalar conditioning frequency is normalized in $\log_{10}$
space with an offset of $6.0$ and a range of $2.778151250383644$. The stored band edges are
$1000000.0$~Hz and $600000000.0$~Hz, which is the 1--600~MHz span quoted throughout this thesis.

\fillin{state why the heatmap uses per-MHz robust statistics while the impedance spectrum uses a
single global log z-score. Both choices are visible in normalization\_stats.json; the reasoning is
not recorded.}

\subsection{Data Splitting, Sampling, and Batch Construction}
\label{subsec:data_splitting_sampling}

The holdout is taken at layout level rather than at row level. Every row belonging to a design ---
that is, all of its anchor frequencies --- is assigned entirely to either the training or the
validation partition, in the ratio 0.9, and the assignment is deterministic given the split seed and
the sorted design keys. The default split seed is 42.

The motivation is recorded in the implementation itself. An earlier version split frequencies within
a design, which placed the same layout on both sides of the holdout and, in the words of that code,
``leaked layout identity into val and is unsuitable for GP / AL residual evaluation''. The change is
dated 2026-07-20 and required a fresh retrain from scratch, with the prior artifacts archived. The
experiment notes give the resulting partition as approximately 22481 training and 2498 validation
layouts; those two figures are stated with a tilde in the notes and are not emitted by any artifact.
\fillin{recompute the exact train/val layout counts for split seed 42 and a 0.9 split, and cite the
emitting artifact rather than the approximate figures in notes.md.}

Batch construction departs from what the configuration file appears to promise. A
$K$-stratification flag is set, but in this code path it only causes the budget of each row to be
precomputed; the weighted sampler is instantiated only when $K$-balancing or frequency-balancing is
requested, and both of those are disabled. Neither stratification nor balancing is therefore active
in exp059.
\fillin{decide whether the intended $K$-stratified sampling was ever exercised, or whether the flag
should be removed; the configuration key and the dataloader behaviour disagree.}

With no balancing sampler active and an epoch budget of 64000 rows, smaller than the training pool,
each epoch draws its rows at random with replacement. Batches additionally carry cross-frequency
pairs: a lookup table maps each training index to alternate-MHz rows of the same design, and the
collate function attaches the paired row to the batch. This pairing is what the cross-frequency
consistency loss of Section~\ref{subsec:objective_physics_reg} consumes.

Loader settings are a batch size of 224 per GPU, a validation batch size of 320, four worker
processes and a prefetch factor of 2, with neither persistent workers nor in-RAM caching.

\section{Multi-modal Variational Autoencoder as a Generative Surrogate}
\label{sec:multimodal_vae_surrogate}

\subsection{Architectural Design}
\label{subsec:architectural_design}

The surrogate is a variational autoencoder \cite{kingma2013auto} whose input modalities are fused by
a Product of Experts \cite{wu2018mvae} and which is conditioned on the decap budget $K$ and the
inspection frequency in the manner of a conditional VAE \cite{sohn2015conditional}. The exp059 model
class, \texttt{MultiInputVAEPoeFreq}, extends the shared multi-input VAE with Product-of-Experts
fusion, adds a frequency expert that writes the private latent dimensions, and removes the U-Net skip
connections.

\paragraph{Structure of the latent space.} The latent dimension is 128, of which 40 dimensions are
private to the heatmap path. Conditioning vectors are 32-dimensional and the frequency encoder uses
16 Fourier bands. The impedance-peak block is not set in the experiment configuration and takes its
constructor default of 8 dimensions. The latent is partitioned as
$z = [z_{\mathrm{shared}} \mid z_{\mathrm{peak}} \mid z_{\mathrm{spatial}}]$ with
\begin{equation}
d_{\mathrm{shared}} = d_{\mathrm{latent}} - d_{\mathrm{hm\text{-}priv}} - d_{\mathrm{imp\text{-}peak}},
\qquad
d_{\mathrm{layout}} = d_{\mathrm{shared}} + d_{\mathrm{imp\text{-}peak}} ,
\label{eq:latent_partition}
\end{equation}
which for these values gives a shared block of 80 dimensions and a layout block of 88.

Decoder input widths follow the partition. The heatmap decoder takes the full latent together with
the concatenated 64-dimensional conditioning vector. The occupancy decoder takes only the 80-
dimensional shared block plus a 32-dimensional conditioning vector. The impedance decoder takes the
88-dimensional layout block $[z_{\mathrm{shared}} \mid z_{\mathrm{peak}}]$, a 32-dimensional
conditioning vector, and a 64-dimensional occupancy graph context. A placement is thus forced through
the shared block alone, spectrum-peak structure is given a private sub-block, and frequency-specific
spatial detail is carried in $z_{\mathrm{spatial}}$, which only the heatmap decoder reads.

\paragraph{Product-of-Experts fusion.} Each modality contributes a Gaussian expert
$\mathcal{N}(\mu_m, \sigma_m^2)$ and the fused posterior takes the standard Product-of-Experts form:
the combined precision is the sum of the expert precisions and the combined mean is the
precision-weighted mean of the expert means \cite{wu2018mvae}. A unit-Gaussian expert
($\mu = 0$, $\log \sigma^2 = 0$) is always appended as an anchor, so the product remains well defined
when only one modality is present. Expert log-variances are clamped to $[-6.0, 1.0]$ and the fused
log-variance to $[-4.0, 2.0]$.
\fillin{state why Product-of-Experts was chosen over mixture-of-experts fusion \cite{shi2019mmvae}
for this problem; the trade-off is discussed in \cite{kutuzova2021defense} but the decision is not
recorded for exp059.}

\paragraph{Modality encoders.} The occupancy expert is a graph encoder over the 52 decap slots laid
out on a $7 \times 8$ grid, with four-neighbour edges, optional self-loops and symmetric degree
normalization, following the graph-convolution formulation of \cite{kipf2016graph}; it is implemented
in plain PyTorch without a graph library. The impedance expert is a chain-graph network over the 231
spectrum bins --- edges between consecutive bins, self-loops, symmetric normalization --- with a
learned per-bin embedding, so neighbouring frequency bins exchange messages in the sense of
\cite{gilmer2017message}. The graph hyperparameters are not set in the experiment configuration; the
applied defaults are a hidden width of 128, three layers and dropout 0.1. The 64-dimensional
occupancy and impedance features are fused before the latent heads by a linear map from 128 to 64
followed by layer normalization and a leaky ReLU of slope 0.1. The heatmap expert is a convolutional
encoder with an $8 \times 8$ spatial bottleneck and dropout 0.05.

\paragraph{Conditioning on budget and frequency.} The budget is conditioned through a discrete
embedding over the 53 possible values $K \in \{0,\dots,52\}$ into the 32-dimensional conditioning
space. The inspection frequency is conditioned through a \texttt{FreqConditioner}: sinusoidal Fourier
features over 16 bands are concatenated with the raw normalized frequency, giving an input of
dimension $1 + 2 \times 16$, and passed through a two-layer SiLU network to 32 dimensions. The vector
handed to the heatmap path is the concatenation of the budget and frequency embeddings, of size 64.
The FiLM and frequency-conditioning modules are imported from an earlier experiment fork (exp043)
rather than defined within exp059, so the exp059 stack is not self-contained in the way the other
forks are.
\fillin{list every cross-experiment import in the exp059 code directory (a grep for
``experiments.exp0'' over that directory settles it); an earlier draft claimed the
frequency-conditioning import was the only one, which is not established.}

Frequency modulation of the decoder uses FiLM \cite{perez2018film}, applied as
\begin{equation}
\mathrm{FiLM}(x) \;=\; x \odot (1 + \gamma) + \beta ,
\label{eq:film}
\end{equation}
where $\gamma$ and $\beta$ are produced from the conditioning vector by a projection whose weight and
bias are zero-initialised except for $\mathrm{bias}[:C] = 1.0$, so that the layer starts close to the
identity. Heatmap FiLM is enabled in the configuration. Multi-scale FiLM is also active, but its
switch does not appear in the experiment configuration and takes a code default of true --- the
architectural change that characterises exp059 is therefore governed by a default rather than by the
experiment's own configuration file. With it enabled, the heatmap decoder is modulated at four
scales: 128 channels at the $8 \times 8$ bottleneck, 64 channels at $16 \times 16$, a 32-channel
stage, and 16 channels at $64 \times 64$.

\paragraph{Decoder path and spatial injection.} The U-Net skip connections were removed, and the
reason is stated in the code: ``No U-Net skips: decoder is latent+FiLM only (matches stage-2 /
generation)''. The corresponding fusion modules were deleted. Spatial information about the placement
instead enters through an occupancy spatial tower, which scatters the 52-dimensional occupancy onto
the $7 \times 8$ grid and fuses occupancy features into the heatmap decoder at $16 \times 16$,
$32 \times 32$ and $64 \times 64$.

\paragraph{Layout-private head.} A layout-private head predicts the heatmap-private latent dimensions
from occupancy, impedance and the conditioning vector, and acts as the distillation student for the
full-encode teacher. Its hidden width is 384, a frequency FiLM is applied to the predicted mean, and
the output is squashed as $\mu_{\mathrm{priv}} \leftarrow 4.0 \tanh(\mu_{\mathrm{priv}} / 4.0)$.

\paragraph{Modality dropout.} During training a Bernoulli keep mask is drawn over the three experts
--- heatmap, occupancy and impedance --- and dropped experts are replaced by the prior. At least one
expert is always kept, and the heatmap expert is force-kept once its protection stage is reached. The
dropout probability is 0.06, annealed from a starting value of 0.06 over 18 epochs, with heatmap
protection from epoch 25 and a reduced value of 0.03 during the heatmap-focus phase. The design
intent is to keep the encoder usable when experts are missing at inference: occupancy is the only
modality available for a placement that has been proposed but not yet simulated, and the
active-learning loop scores such placements through the occupancy-only encode path.
\fillin{cite an ablation or per-path evaluation showing that the occupancy-only encode path degrades
without modality dropout. No such artifact exists, so the enabling role of this term is an untested
design assumption rather than a measured effect.}

\paragraph{Discrete occupancy read-out.} Occupancy is decoded to a hard top-$K$ binary vector before
the heatmap decode, so that the spatial prediction is conditioned on a CAD-realisable placement
rather than on a soft one. The read-out is fully vectorised, using a descending-sort threshold that
supports per-row variable $K$ and avoids Python loops and device synchronisation on the decode path.
A straight-through mode is available, in which the forward value is the hard top-$K$ vector and the
gradient flows through the soft occupancy \cite{bengio2013straight}; in exp059 training that mode is
off.

\paragraph{Encode paths available at inference.} The model exposes several encode paths. Four matter
here: encoding from all modalities; cross-modal encoding from the heatmap, the impedance, or
occupancy and impedance together; encoding of the layout latent from occupancy and impedance with the
private head; and \texttt{encode\_occupancy\_latent}, which encodes from occupancy alone. The last is
the path on which the active-learning loop and any forward design query depend.

\figplace{Block diagram: three modality experts, PoE fusion with prior anchor, structured latent split, FiLM-conditioned heatmap decoder with occupancy spatial tower.}

\subsection{Objective Functions and Physics-Informed Regularization}
\label{subsec:objective_physics_reg}

\paragraph{Reconstruction terms.} The reconstruction objective is a weighted sum over the three
modalities plus an occupancy consistency term,
\begin{equation}
\mathcal{L}_{\mathrm{recon}}
= w_{\mathrm{hm}} \mathcal{L}_{\mathrm{hm}}
+ w_{\mathrm{occ}} \mathcal{L}_{\mathrm{occ}}
+ w_{\mathrm{imp}} \mathcal{L}_{\mathrm{imp}}
+ w_{\mathrm{occK}} \mathcal{L}_{\mathrm{occK}} ,
\label{eq:recon}
\end{equation}
with heatmap weight 5.0, occupancy weight 7.0, impedance weight 4.0 and occupancy-$K$ consistency
weight 2.0. A heatmap-focus phase beginning at epoch 25 overrides these to heatmap 7.0, impedance
4.0, cross-frequency 1.1 and dynamic range 3.0.
\fillin{give the design rationale for the numeric loss weights (5.0 / 7.0 / 4.0 / 2.0 and the
focus-phase overrides). The experiment notes explain the capacity and FiLM changes of exp059 but not
how these weights were arrived at; a weight-sweep artifact would settle it.}

Occupancy reconstruction uses a focal binary cross-entropy with focusing parameter 1.5 and a
positive-class weight $(52-K)/K$ clamped to $[0.5, 10.0]$, which compensates for the strong class
imbalance of a small-$K$ binary target. The $K$-consistency term is a hinge that penalises
$\mathrm{relu}(0.5 - \ell)$ on target-on slots and $\mathrm{relu}(\ell + 0.5)$ on target-off slots,
where $\ell$ is the decoded occupancy logit.

The impedance loss is a composite over the 231 bins. Its components are a derivative term of weight
1.4; a top-$K$ term over the 20 largest bins with weight 4.5; an under-prediction penalty of 2.8; a
concavity term of 2.5; a frequency-weighting exponent of 2.0; a dual top-$K$ term of 0.75; and
peak-index and peak-magnitude terms of 1.5 and 2.0 over the 8 strongest peaks. The peak terms are
ramped in rather than applied from the start, beginning at epoch 25 and reaching full weight over 22
epochs, with the peak-focus phase also starting at epoch 25.

Heatmap reconstruction adds several shape-sensitive terms to the pixel loss: a Pearson-correlation
term of weight 0.75, a gradient-vector term of 2.0 and a gradient-direction term of 1.0, the latter
two with a minimum gradient magnitude of 0.08 and a Huber delta of 0.5. Extreme percentiles are
treated separately by top-region Huber terms at the 0.95 and 0.05 quantiles, each of weight 2.5 and
Huber delta 0.5. A soft-argmax peak-location loss carries weight 2.5 at temperature 0.035; the
corresponding valley-location loss is disabled with weight 0.0.

\paragraph{Latent regularization.} The KL divergence is computed per dimension as
\begin{equation}
\mathrm{KL}_j = -\tfrac{1}{2}\left(1 + \log \sigma_j^2 - \mu_j^2 - \exp(\log \sigma_j^2)\right),
\label{eq:kl}
\end{equation}
clamped from below by a free-bits floor of 0.12 before averaging, so that a dimension carrying less
than that much information is not pushed further toward the prior. Beta annealing is disabled and
$\beta$ is fixed at 0.03, giving a fixed $\beta$-weighted KL in the sense of \cite{higgins2017beta}.
Three further latent regularizers enter the total: a hinge
$w_{\mathrm{hinge}} \, \mathrm{relu}(|\mu| - \tau)^2$ with weight 0.1 and threshold 4.0; a batch-mean
bias term $w_{\mathrm{bias}} \, \overline{\mu}^2$ with weight 0.05; and a per-expert KL averaged over
experts with weight 0.02. A sigma regularizer of weight 1.5 keeps the fused posterior standard
deviation inside a band around a target of 0.45, penalising values below the target and above
$1.3 \times$ the target.
\fillin{explain the choice of sigma regularizer target 0.45 and free-bits floor 0.12; no artifact
records how they were selected.}

\paragraph{Cross-frequency consistency, distillation and jitter.} Cross-frequency training decodes
the latent at an alternate MHz value and compares the result with the paired alternate-MHz heatmap
supplied by the pair lookup of Section~\ref{subsec:data_splitting_sampling}. Its weight is 1.2 from
epoch 20, with a layout-mixing probability of 0.45, and it re-uses the encoded latent. Distillation
transfers the full-encode teacher onto the layout student with a latent term of weight 2.5 --- the
private dimensions multiplied by 3.0 and the log-variance component by 0.3 --- and an output term of
weight 2.5, implemented as a Huber loss between the layout decode and the detached teacher decode.
Frequency jitter is applied on the decode side with probability 0.6 and a $\log_{10}$ standard
deviation of 0.1.

\paragraph{Spectral anti-blur term, inactive in this experiment.} A two-dimensional FFT
log-magnitude $L_1$ anti-blur loss is implemented, together with an optional Gaussian bump that would
emphasise a mid-band frequency window. It does not act in exp059. Its weight key is absent from the
experiment configuration, so the weight lookup returns its default of 0.0, and the function returns
nothing for any non-positive weight; the term never enters the total loss. This has to be stated,
because the experiment notes present the frequency-weighted spectral loss as one of the three
headline changes of exp059 and quote a weight of 1.5, a mid-band boost of 1.5, a centre of 200~MHz
and a width of 80~MHz. Those keys do not exist in the configuration as it stands.
\fillin{reconcile the experiment notes, which list the spectral anti-blur loss as a headline exp059
change, with the configuration, in which the weight key is absent and the term inactive: either
restore the keys and rerun, or correct the note.}

\paragraph{Physics regularizers and their absence from the objective.} Three physics terms are
implemented in a dedicated loss module. \emph{Radius influence} penalises reconstructed foreground
pixels that are cold where a frequency-conditioned critic judges decap coverage to be weak.
\emph{Critic supervision} trains that critic on the inverted min--max normalised foreground of the
target heatmap; the critic is a dual-branch convolutional network taking the occupancy grid together
with a 16-dimensional frequency embedding and producing a spatial effectiveness map, and it is
detached inside the radius-influence term so that VAE gradients do not corrupt it.
\emph{Anti-resonance} encodes the requirement that every impedance peak be preceded in frequency by a
resonance dip:
\begin{equation}
\begin{aligned}
s^{\mathrm{peak}} &= \min\!\left(\mathrm{relu}(c - l), \; \mathrm{relu}(c - r)\right), \qquad
s^{\mathrm{dip}} = \min\!\left(\mathrm{relu}(l - c), \; \mathrm{relu}(r - c)\right), \\
\mathcal{L}_{\mathrm{ar}} &= \mathrm{mean}\!\left(s^{\mathrm{peak}} \cdot \exp(-10.0 \cdot d^{\mathrm{prior}})\right),
\end{aligned}
\label{eq:anti_resonance}
\end{equation}
where $c$, $l$ and $r$ are the centre, left and right bin values of the impedance channel and
$d^{\mathrm{prior}}$ is the running cumulative maximum of prior dip scores.

The experiment configuration assigns non-zero weights to all three: 1.0 for radius influence, 2.0 for
critic supervision and 0.5 for anti-resonance, with a critic warm-up of 11 epochs, a slope anneal
over 22 epochs and a foreground clip at $-1.6791870594024658$. The physics module is instantiated,
its parameters reach the optimizer, and they are stored in the checkpoints. None of the three terms
nevertheless contributes to the training loss in this experiment. The exp059 trainer replaces the
epoch function of the shared training core with a variant whose signature no longer carries the
physics arguments and which calls the loss with the physics module unset; the loss guards the
radius-influence, critic-supervision and anti-resonance terms on the presence of that module, so none
of them is added.

The exp059 surrogate is consequently not physics-informed in its training objective, and this thesis
does not describe it as such. The physics prior that is genuinely active in an objective sits in
Stage 2, which as implemented runs on the legacy \texttt{exp038\_true\_multi} stack rather than on
exp059 (Section~\ref{subsec:opt_problem_formulation}). There, the latent optimizer applies an
anti-resonance term of weight 0.5, overridable by the checkpoint configuration; see
Section~\ref{subsec:opt_procedure_feasibility}.
\fillin{state whether disabling the physics terms in the exp059 training loop was a deliberate
ablation or an unintended consequence of the epoch-function replacement. Nothing in the repository
explains the disconnect, and the answer changes how this subsection should be framed.}

\subsection{Training, Inference}
\label{subsec:training_inference}

\paragraph{Configuration and code organisation.} Each experiment reads a configuration file from its
own directory, and an environment override may be layered on top of it at runtime. That override is
the mechanism by which the active-learning fine-tune injects its own settings
(Section~\ref{sec:active_learning}). The exp059 trainer specialises the shared training core by
replacing several of its components: the configuration object, the model builder, the epoch function,
the heatmap loss, the total loss, the statistics hook, the dataloader factory, the configuration
loader and the impedance loss.

\paragraph{Optimizer and schedule.} The optimizer is AdamW with weight decay $1 \times 10^{-4}$,
selected as fused where available, then in its \texttt{foreach} form, and otherwise as a plain
implementation. The scheduler reduces the learning rate on plateau in minimising mode, with a factor
of 0.5, a patience of 15 epochs and a floor of $1 \times 10^{-5}$. A separate impedance learning rate
can be enabled after the peak-focus epoch; it is enabled, with a multiplier of 1.0 and a peak-focus
epoch of 25.

The configuration as it stands specifies a learning rate of $1.5 \times 10^{-5}$, 510 epochs, a
learning-rate reset on resume, and a resume checkpoint pointing at the experiment's last model. The
experiment notes state that the exp059 training run used a learning rate of $2 \times 10^{-4}$, 500
epochs and no resume checkpoint. The two no longer agree.
\fillin{state which learning rate, epoch count and resume setting produced the checkpoint whose
results are reported --- the notes give 2e-4 / 500 / null and the configuration gives 1.5e-05 / 510 /
last\_model.pt --- and cite the run log that settles it.}

\paragraph{Numerical stability.} The policy is deliberately conservative. Mixed precision, TF32 and
graph compilation are all off, training is forced to fp32, non-finite batches are skipped, and
gradients are clipped to a norm of 1.0. Log-variances are clamped to $[-10.0, 10.0]$, the
reconstruction latent has a soft floor of $-8.0$, and the loss is capped at $1000000.0$ for
finiteness.
\fillin{state why mixed precision and TF32 were disabled for exp059 --- whether a specific
instability was observed, and in which run.}

\paragraph{Curriculum.} The milestone epochs in the current configuration are: cross-frequency
consistency from epoch 20; heatmap focus, heatmap protection under modality dropout, the impedance
peak terms and the end of beta scheduling from epoch 25; and warm-ups of 11 epochs for the physics
critic and for the penalty terms.

\paragraph{Validation and monitoring.} Validation runs on the first epoch, on each checkpoint
interval and on the final epoch, with a checkpoint interval of 5 epochs and no checkpoint pruning.
The off-anchor evaluation runs every 5 epochs at 155.0, 250.0 and 265.0~MHz with weights 2.5, 3.0 and
2.5 respectively, over at most 20 batches, and evaluates the occupancy-only decode path. The name is
only partly accurate for this frequency list: 155.0 and 265.0 are absent from the 24 training anchors
of Section~\ref{subsec:multifreq_dataset} and are genuinely held out, whereas 250.0 is a training
anchor and carries the largest of the three weights, so it acts as an in-distribution reference
inside the same metric. The experiment directory also retains a best-off-anchor checkpoint, so this
metric doubles as a checkpoint-selection criterion rather than serving only as a monitor.
\fillin{state why 250~MHz, a training anchor, sits in the off-anchor evaluation set with the highest
weight of the three, and confirm exactly which aggregate quantity best\_off\_anchor\_model.pt is
selected on.}
Each run writes per-epoch records of the total and component losses, the heatmap peak split, the
impedance split, the off-anchor evaluation and per-epoch timing into the experiment's metrics
directory.

\paragraph{Distributed training.} Distributed data-parallel training is enabled with a base batch
size of 160 and linear learning-rate scaling, so the learning rate is multiplied by
$(\text{batch size} \times \text{world size}) / 160$. Initialisation uses the NCCL backend with a
default 30-minute timeout, and the model is wrapped with unused-parameter detection enabled.
\fillin{state why unused-parameter detection was enabled in the distributed wrapper. The setting
carries no comment, note or commit message giving a reason, so the natural explanation --- expert
branches left unused by modality dropout --- is not evidenced.}

\paragraph{Encode-path mixing during training.} The layout-mixing probability is 0.65 and the
occupancy-only encode probability is 0.4, and this chapter takes these and the neighbouring keys from
the configuration file. The experiment notes and the project handover document both describe the
exp059 capacity phase as full-encode only, with both probabilities at zero.
\fillin{state which encode-path mixing setting produced the reported capacity results. The
configuration gives layout\_train\_prob 0.65, occ\_only\_encode\_prob 0.4, occupancy-only evaluation
enabled and distillation weights of 2.5, whereas the notes and the handover document describe all of
these as zero or disabled for the capacity phase.}

\fillin{document the training hardware and software environment: GPU model and count, driver, CUDA
and PyTorch versions, and wall-clock time per epoch. No committed artifact records this; the timing
CSV gives per-epoch durations but not the machine.}

\fillin{report the exp059 model parameter count. It is printed at runtime by the training core but
not stored in any artifact.}

\subsection{Auxiliary Guidance Network for Latent-Space Optimization}
\label{subsec:aux_guidance_network}

No module, class or function of this name exists in the current model track; the term appears only in
the thesis outline. The nearest implemented component is a feed-forward impedance surrogate that maps
an occupancy vector directly to an impedance spectrum, and which Stage 2 uses to score candidate
placements without invoking the VAE impedance head. It exists only in the legacy experiment forks
exp037, exp038, exp039, exp040 and exp041, and is absent from exp057 through exp060.

In the exp038 fork it maps occupancy $(52,)$ to an impedance log-$z$ spectrum $(1, 231)$ and does not
use the inspection frequency. Its hidden widths are $(256, 512, 512, 512, 256)$ with dropout 0.1, and
it is trained for 200 epochs at batch size 128 with a learning rate of $3 \times 10^{-4}$, plateau
patience 15, factor 0.5 and floor $1 \times 10^{-5}$, weight decay $1 \times 10^{-4}$, a top-$K$ term
over 20 bins with weight 7.0, and an under-prediction penalty of 2.8. Stage 2 loads its checkpoint
from the exp038 experiment directory whenever the surrogate is enabled, which is the default, and
otherwise falls back to the VAE impedance decoder head.

\fillin{decide what this subsection is to contain. Either (a) rename the heading to ``Impedance
Surrogate for Latent-Space Optimization'' and present the exp038 surrogate above as the component,
stating that it belongs to the legacy Stage-2 stack and not to exp059; or (b) implement and describe
an exp059-native guidance network. As it stands there is no exp059 component under this name.}

\section{Latent-Space Optimization of Decoupling Capacitor Placement}
\label{sec:latent_opt_decap}

\subsection{Problem Formulation}
\label{subsec:opt_problem_formulation}

Stage 2 inverts the surrogate. With all network weights frozen, the search variable is the latent
vector $z$ rather than the placement itself, following the latent-space inverse-design template of
\cite{gomezbombarelli2018automatic}. For a fixed budget $K$, the decoded occupancy probabilities
$p(z) \in [0,1]^{52}$ are read out to a placement by a hard top-$K$ selection,
\begin{equation}
\mathbf{b}(z, K) = \mathrm{top}_K\!\left(p(z)\right) \in \{0,1\}^{52},
\qquad \sum_{i} b_i = K ,
\label{eq:topk_readout}
\end{equation}
and the predicted spectrum $\hat{Z}(f)$ over the 231 bins is compared with the target mask
$Z_{\mathrm{target}}(f)$. A candidate is declared feasible only when the predicted curve lies at or
below the feasibility boundary at \emph{every} one of the 231 bins,
\begin{equation}
\hat{Z}(f_j) \le Z_{\mathrm{boundary}}(f_j) \quad \text{for all } j = 1,\dots,231,
\qquad
Z_{\mathrm{boundary}} = Z_{\mathrm{target}} - m,
\label{eq:feasibility_design}
\end{equation}
with a boundary margin $m = 0.1$. The objective maximises the margin below the boundary, and the loss
is evaluated in the normalized log space, so the base term rewards margin and penalises excursions
above the boundary,
\begin{equation}
\mathcal{L}_{\mathrm{base}}
= -\,w_{\mathrm{gap}} \cdot \mathrm{gap}
+ w_{\mathrm{exceed}} \cdot \mathrm{exceed},
\qquad
w_{\mathrm{gap}} = 1.0, \; w_{\mathrm{exceed}} = 25.0 ,
\label{eq:stage2_base}
\end{equation}
where the exceed component is a mean of $\mathrm{relu}(\hat{Z} - Z_{\mathrm{boundary}})$ raised to
the power 2.0 over the bins, plus a maximum-excursion term.
\fillin{write out the exact form of the maximum-excursion component of the exceed term and of the gap
term; the extraction recorded the weights and the mean component but not the full expression.}

The target mask is loaded from a stored array, is required to have shape $(231,)$ after squeezing, is
interpreted in ohms, and is converted into the model's normalized log space by
$(\log Z_{\mathrm{target}} - \mu_{\log}) / \sigma_{\log}$. The stored array has shape $(231, 1)$ and
dtype float64, with minimum 0.8 and maximum 49.99999999999999.
\fillin{document the provenance and definition of the target impedance mask: who set the 0.8 ohm
floor and the 49.99999999999999 ohm ceiling, against which specification, and with which script. The
.npy file carries no unit metadata and no generating script was found.}

\paragraph{Scope limitation.} Stage 2 as implemented targets the legacy
\texttt{exp038\_true\_multi} stack: it defaults to that experiment for both the VAE occupancy decoder
and the impedance surrogate checkpoint, and imports its VAE class from that fork. It is not wired to
exp059, a limitation also recorded in the project's latent-optimization documentation. Everything
described in this section is therefore the method as implemented against exp038, and no end-to-end
inverse-design result may be attributed to the exp059 surrogate. A second, smaller inconsistency
points the same way: the Stage-2 default normalization statistics path names a dataset directory that
does not exist in the current checkout, whereas the exp059 training set is the unbounded set of
Section~\ref{subsec:multifreq_dataset}.
\fillin{record the work needed to wire Stage 2 to the exp059 checkpoint, including the correct
normalization statistics path; the Stage-2 default points at datasets/data\_multifreq\_norm/, which is
absent, while training uses datasets/data\_multifreq\_train\_norm\_unbounded.}

\subsection{Optimization Procedure and Feasibility Assessment}
\label{subsec:opt_procedure_feasibility}

The search is gradient descent on $z$ with everything else frozen. Budgets $K = 1 \dots 25$ are
searched, each with 32 candidate seeds optimized as a single batch over 1200 steps at a learning rate
of $5 \times 10^{-2}$, with gradients clipped at 5.0 and an initial shared-latent temperature of 0.8.
Seeds may be drawn randomly per run and recorded in the run configuration, or fixed to $0 \dots N-1$;
the random mode is the default.
\fillin{reconcile the three budget ranges in the framework: Stage 2 searches $K = 1 \dots 25$, active
learning acquires $K = 3 \dots 10$, and the dataset is filtered to $K \le 30$. Nothing in the
repository explains the relationship.}

\paragraph{Straight-through top-$K$ read-out.} The read-out of
Equation~\ref{eq:topk_readout} is not differentiable, and this is the load-bearing detail of Stage 2.
The implementation uses a straight-through estimator \cite{bengio2013straight},
\begin{equation}
\mathbf{b} = \mathrm{top}_K^{\mathrm{hard}}(p) + p - \mathrm{sg}[p],
\label{eq:ste_design}
\end{equation}
where $\mathrm{sg}[\cdot]$ denotes the stop-gradient operator: the forward value is the hard top-$K$
vector and the gradient is passed to $p$ unchanged. The reason is stated in the code --- ``Without
this the top-K argsort severs the gradient and $z$ never moves on the spectrum target.'' Relaxing the
discrete selection with a Gumbel-softmax or concrete relaxation \cite{jang2016gumbel,
maddison2016concrete} is an alternative that is not used here.

One qualification matters for interpreting Stage 2. The estimator is applied only on the surrogate
branch: when the impedance surrogate is loaded, the forward pass evaluates the surrogate on the
straight-through top-$K$ occupancy, whereas without it the impedance comes from the VAE decoder head
and the occupancy read-out is not in the gradient path at all. The surrogate is enabled by default,
so the straight-through branch is the default, but the fallback branch does not exercise it.

\paragraph{Regularization of the search.} Several terms keep the search inside the region where the
decoder is trustworthy and keep the 32 seeds from collapsing onto one another. A latent prior term of
weight 0.05 penalises $((z - \mu)/\sigma)^2$ against the aggregate posterior rather than against the
unit prior, and a posterior-boundary hinge of weight 10.0 activates beyond a Mahalanobis-style radius
of 2.5. Spectrum-shape terms comprise a shape regularizer of weight 2.0 with a minimum log-spectrum
standard deviation of 0.12, which penalises degenerate flat spectra, and a tracking term of weight
2.0 for mean-squared alignment with the target. Peak alignment is enforced by a peak loss of weight
1.0, a peak-index term of 3.0, a peak-magnitude term of 2.0 and a dual top-$K$ term of 2.5, over the
8 strongest peaks and the 20 largest bins; all of these may be overridden by environment variables.
Seed diversity is encouraged by penalising the mean pairwise cosine similarity across the parallel
seeds with weight 0.35, and occupancy confidence with weight 0.5. Finally the anti-resonance physics
prior of Equation~\ref{eq:anti_resonance} is applied with weight 0.5, overridable by the checkpoint
configuration. This is the one physics term that is active in an optimization objective in this work,
and it belongs to Stage 2 as implemented against exp038; as
Section~\ref{subsec:objective_physics_reg} records, no physics term enters the exp059 training
objective.

\paragraph{Selection and reporting of solutions.} Among the feasible candidates for a budget, the
selection metric is the maximum predicted impedance in ohms across the spectrum: candidates are
ranked by their peak and the lowest peak wins. Ranking by the mean is available as an alternative. If
no candidate for a budget is feasible, the run writes an explicit no-solution record for that budget
rather than emitting a best-effort placement, which keeps infeasibility distinguishable from a poor
solution in the downstream reports.

Surrogate feasibility is a screening verdict and not a sign-off. A placement that Stage 2 declares
feasible has been checked only against the surrogate's own prediction, and requires ground-truth
re-simulation in ECADSTAR before any claim about the physical board is made.

\figplace{Stage-2 loop: frozen decoder, gradient on z, straight-through top-K read-out into the impedance surrogate, feasibility test against the target mask.}

\section{Experimental Setup}
\label{sec:experimental_setup}

The model track used for all training and active learning reported in this thesis is
\texttt{exp059\_capacity\_freq}, referred to as exp059 throughout.
\texttt{exp057\_structured\_graph} and \texttt{exp058\_asymmetric\_kl} are legacy, and
\texttt{exp060\_multitype\_occ} is an exploratory fork. Each experiment directory is a self-contained
copy of the training stack, so a change in one fork does not propagate to the others, and checkpoints
are not cross-loadable between them.

Training data is the unbounded normalized set described in Section~\ref{subsec:multifreq_dataset}:
582997 manifest rows over 24979 layouts and 24 anchor frequencies, split by layout in the ratio 0.9
with seed 42. Two configuration files govern the two phases of the experiment, \texttt{config.yaml}
for the capacity phase and \texttt{config\_al\_finetune.yaml} for the active-learning fine-tune, the
latter injected at runtime through the configuration-path override. A checkpoint trained under the
capacity configuration alone is not directly usable for active-learning scoring, which needs the
occupancy-only encode path.

\fillin{state the compute environment used for the reported runs: GPU model and count, host memory,
CUDA and PyTorch versions, and total training wall-clock time. This is not recorded in any committed
artifact and must not be taken from a live shell.}

Simulation of newly proposed layouts runs on ECADSTAR under Windows, driven from WSL. The
active-learning configuration names the design file H-shape.erf, the powerbus Power\_GND and the port
IC1\_Port1, requests heatmap-only PEB export, and allows a simulation wait of 18000~s. The simulation
steps of the pipeline cannot be executed without that host.

The residual Gaussian process used for acquisition (Section~\ref{sec:active_learning}) is implemented
with GPyTorch \cite{gardner2018gpytorch} as an approximate GP trained with a variational evidence
lower bound.

All numeric settings quoted in this chapter are taken from the experiment configurations, the dataset
metadata and the active-learning configurations. Several of the project's internal reference
documents describe earlier experiment generations and disagree with those artifacts on anchor sets,
loss weights and acquisition budgets; they are used for method description only and not as a source
of numbers.

\section{Active Learning for scaling the model}
\label{sec:active_learning}

The surrogate is trained on a fixed pool of simulated layouts, and the layouts that matter for
inverse design --- the placements an optimizer wants to propose --- are not necessarily well
represented in that pool. The active-learning loop closes this gap in the standard pool-based way
\cite{settles2009active, batchal2022batch}: generate a large pool of candidate layouts, score them,
send only the highest-scoring batch to the expensive ECADSTAR simulator, ingest the labels, and
fine-tune the surrogate on the enlarged set.

\paragraph{Cycle structure and control.} A cycle is described here as eight steps: per-$K$
acquisition, which covers candidate generation, inference and selection; PEB construction; ECADSTAR
simulation; ingest with normalization and pre-fine-tune evaluation; overlay construction; fine-tune;
and post-fine-tune evaluation with the decision report. The step count is computed at runtime and is
eight only when fine-tuning and post-fine-tune evaluation are both enabled, and seven or six
otherwise. A propose-only mode runs generation and inference alone, with no simulation and no
fine-tune, and is the mode used for dry runs. Three configurations are provided: a residual-GP
acquisition (the default and primary run), a Monte-Carlo acquisition, and an equal-budget random
control.

\paragraph{Candidate pool.} For exp059 the pool spans budgets $K = 3$ to $K = 10$, sweeping all
budgets in one pool, with 1600 candidates per budget, 80 selected per budget and a total ECAD budget
of 640 simulations per cycle. Candidates are random $K$-hot occupancy vectors over the 52 slots.
Generation is seeded deterministically per iteration and per budget as
$\texttt{seed} = 42 + \texttt{iteration} \times 10000 + K$, so a cycle can be reproduced exactly.
Each candidate is also assigned an inspection frequency from the grid
$[120, 150, 155, 170, 180, 200, 230, 250, 265, 270, 280, 300]$~MHz under per-$K$ strata quotas
$120{:}30$, $150{:}50$, $155{:}100$, $170{:}40$, $180{:}50$, $200{:}90$, $230{:}90$, $250{:}140$,
$265{:}130$, $270{:}80$, $280{:}70$, $300{:}70$. An exploration slice adds 160 further candidates per
budget, drawn from a continuous 120--320~MHz range quantized to 5~MHz with band edges
$[120, 180, 230, 280, 320]$. Selection applies the frequency quotas together with a minimum
uncertainty percentile of 40, and where pool strata are given without matching ECAD-side strata, the
latter are auto-scaled from the pool quotas to that budget's selection count.
\fillin{give the rationale for the frequency strata quotas --- why 250~MHz receives 140 candidates
per $K$ and 155~MHz receives 100. No artifact states it.}

These budgets differ from the written active-learning documentation, which is exp057-era and states
400 candidates per budget, 10 selected per budget, 80 ECAD simulations per cycle and off-anchor
frequencies of 90, 265, 435 and 510~MHz. The values above are those of the exp059 configuration.
\fillin{update docs/active-learning.md, whose acquisition budgets and off-anchor frequencies predate
the exp059 configuration used here.}

\paragraph{Acquisition by a residual Gaussian process.} The primary exp059 run acquires by predicted
error. The acquisition dispatch supports three modes --- residual GP, random and Monte-Carlo --- and
defaults to Monte-Carlo when unset. The residual-GP mode does not ask the VAE how confident it is: it
fits a separate Gaussian process whose regression target is the surrogate's own error on labelled
data, and ranks candidates by that predicted error. The configured settings use the predicted-heatmap
feature path, the 99th-percentile absolute error over the physical foreground as target, a sparse
variational GP regressor, $\kappa = 0.5$, 4000 fitting points, 256 inducing points, 80 fitting
epochs, seed 0, a novelty weight of 0.0, and a score mode of posterior mean. Other residual targets
are available in the implementation: peak-pixel error in three variants, a structural target, mean
absolute and mean squared error, and impedance mean squared error.

The acquisition function implemented is the upper confidence bound
\begin{equation}
a(x) = \mu_e(x) + \kappa \, \sigma_e(x),
\label{eq:gp_ucb}
\end{equation}
in the form of \cite{srinivas2010gp}, where $\mu_e$ and $\sigma_e$ are the posterior mean and
standard deviation of the error GP. Under the configured score mode the acquisition returns $\mu_e$
alone, and the novelty term enters only under the combined UCB-plus-novelty mode. The production
exp059 run therefore ranks purely by predicted error and does not use $\sigma_e$. One implementation
docstring describes the mode as ranking by the full upper confidence bound, which does not match the
resolved behaviour.
\fillin{correct the inference-pool docstring, which states that residual-GP candidates are ranked by
$\mu_e + \kappa\,\sigma_e$, whereas the configured score mode resolves to the mean alone with sigma
unused; docs/gp-error-surrogate.md states the mean default explicitly.}

The GP inputs come from a feature path that requires no simulation: the model encodes the candidate
occupancy through the occupancy-only path, decodes it, re-encodes the model's \emph{own} predicted
heatmap, and uses the resulting posterior mean as the GP input. This is what allows candidates to be
scored before any ECAD label exists. The fitted GP is cleared at the start of every per-$K$
acquisition, so each cycle performs a fresh fit rather than re-using a cached one.
\fillin{confirm the reason for clearing the cached GP fit each cycle --- the obvious candidate is
that fine-tuning shifts the latent space in which the GP is fitted, but only the clearing itself is
recorded, with no comment or note stating why.}

This construction is assembled from three separate literatures rather than taken from one canonical
precedent. The sparse variational GP that makes it tractable at candidate-pool scale comes from the
inducing-point variational formulation of \cite{titsias2009variational} and its stochastic extension
\cite{hensman2013bigdata}, implemented here in GPyTorch \cite{gardner2018gpytorch}. The acquisition
rule is GP-UCB \cite{srinivas2010gp}. Modelling a separate discrepancy term rather than trusting the
primary model's own uncertainty is the model-discrepancy decomposition of
\cite{kennedy2001calibration}. No source found in the literature review combines them as a unit for
active learning of a neural surrogate in engineering design, and this chapter does not claim one.
Modelling error with a second model is also a departure from asking the network itself, for example
by Monte-Carlo dropout \cite{gal2016dropout}, which is the mechanism the Monte-Carlo mode implements.
\fillin{state the design rationale for choosing the residual GP over MC-dropout uncertainty in the
author's own words; the framework document motivates it, but the decision is not recorded as a
comparison in any artifact.}

\paragraph{Controls.} The random control assigns uniform random scores from a generator seeded by the
run seed plus the iteration index, preserving the ECAD budget exactly so that the comparison is
equal-budget. The Monte-Carlo mode retains its parameters in the exp059 configuration: 4 forward
passes, 4 latent passes and 8 decoder passes with decoder dropout enabled, an automatic
worst-candidate score metric, and score weights of 1.0 for decoder-dropout reconstruction error, 1.0
for latent heatmap mean squared error, 1.0 for the variance of the 99th percentile, 0.5 for its
range, 0.5 for peak-location spread, 0.25 for the peak coefficient of variation and 0.35 for prior
tension.

No comparative claim is made here about the three acquisition modes. The exp059 acquisition A/B
remains an open task and the project's claim ledger records the GP-versus-Monte-Carlo comparison as
uncertain, so this chapter describes the mechanisms and leaves their relative merit to be
established.

\paragraph{Scoring configuration.} Scoring runs in layout inference mode at a reference frequency of
200.0~MHz, with a shared-latent temperature of 1.5 and a heatmap grid of 64, using a $64 \times 64$
boolean board mask and the normalization statistics of the unbounded training set, so that scoring
uses the same normalization the model was trained under.

\paragraph{Simulation, ingest and fine-tune.} Selected layouts are exported as PEB batches and
simulated by ECADSTAR, as described in Section~\ref{sec:experimental_setup}. Labels are ingested and
robustly normalized, then merged into an overlay dataset
(\texttt{datasets/data\_multifreq\_al\_overlay\_exp059}) which is mixed into training with a
per-sample weight of 128.0 on the overlay rows.
\fillin{state the reason for the 128.0 overlay sample weight and how the value was chosen; both the
AL configuration and the fine-tune configuration set it, but no artifact records the intent or a
comparison against the training code's own default.}
The fine-tune resumes from the last checkpoint and trains for 45 epochs beyond the checkpoint's own
epoch count. The merged runtime configuration is written out and passed to the trainer through the
configuration-path override.

The fine-tune's encode-path mixing probabilities differ between the two configurations that define
them. The experiment's fine-tune configuration sets a layout-mixing probability of 0.9 and an
occupancy-only encode probability of 0.7, together with latent and output distillation weights of
2.5, occupancy-only evaluation, and early stopping with patience 3, minimum delta 0.01 on the
layout-cross criterion. The active-learning configuration sets 0.65 and 0.4, and the runtime merge
copies the active-learning values over the experiment file, so the executed values are 0.65 and 0.4.
\fillin{state whether the executed fine-tune mixing values of 0.65 / 0.4 were intended; the
experiment's fine-tune configuration, the project handover document and CLAUDE.md all give 0.9 / 0.7,
and the runtime merge silently overrides them.}

\paragraph{Reporting.} Decision reporting is gated on minimum sample sizes of 30 ranked candidates
and 20 ECAD labels, and re-scores only the newly ECAD-labelled layouts. Every cycle writes
\texttt{scored\_candidates.json}, \texttt{acquisition\_rank\_quality.json}, pre- and post-fine-tune
off-anchor evaluations, a cycle summary, a cycle evaluation report, a decision report and a decision
ledger. The results chapter draws its active-learning numbers from these files.

The loop has been executed on exp059. The run \texttt{al\_exp059\_gp\_error\_001} contains four
iterations, and the summary of the fourth records residual-GP acquisition over budgets
$[3,4,5,6,7,8,9,10]$ with 1600 candidates and 80 selections per budget, 12800 candidates in total and
640 selected. This is stated only as evidence that the pipeline described above runs end to end; no
performance figure from those runs belongs in this chapter.

\figplace{Active-learning cycle: candidate generation, residual-GP scoring, PEB export, ECADSTAR simulation, overlay build, fine-tune, decision report.}

% ---------------------------------------------------------------
% NOTE ON TRACEABILITY: all file:line references formerly inline in this chapter were moved to
% the Implementation Index appendix. Source of record:
% scratchpad/thesis_phase1/ch3-impl-index.json (chapter "ch3"). If a claim in this chapter is
% challenged, the artifact and line range are in that index under the same section heading.
%
% FILL IN checklist:
%   - Board description (H-shape test vehicle): geometry, stackup, port definition. Needs a
%     committed board-spec artifact or the author's own description.
%   - ECADSTAR bulk data-generation setup (solver, excitation, meshing, heatmap export). Only the
%     AL-side CLI invocation is documented (docs/data-pipeline.md sec.5).
%   - Stored heatmap tensor resolution as a dataset property (dataset_meta.json has no shape field).
%   - Whether the 80 MHz anchor shortfall (8480 vs 24979) is deliberate.
%   - Anchor-count conflict: dataset_meta.json (24) vs docs/dataset.md sec.4 (16) vs
%     docs/normalization-and-losses.md sec.1.4 (six anchors).
%   - Why heatmaps use per-MHz robust stats but impedance uses one global log z-score.
%   - Exact train/val layout counts for seed 42, train_split 0.9 (notes.md gives ~22481/2498).
%   - stratify_by_k conflict: config.yaml:18 true vs dataloader_base.py:301,353-362 (no sampler).
%   - Complete list of cross-experiment imports in experiments/exp059_capacity_freq/codes/ (grep for
%     "experiments.exp0"). The earlier "only one cross-experiment import" claim was withdrawn.
%   - Ablation or per-path evaluation showing the occupancy-only encode path degrades without
%     modality dropout. No such artifact exists; the enabling role is an untested assumption.
%   - Rationale for PoE over MoE fusion for this problem.
%   - Rationale for the numeric loss weights 5.0 / 7.0 / 4.0 / 2.0 and the focus-phase overrides.
%   - Rationale for sigma_reg_target 0.45 and free_bits 0.12.
%   - Spectral anti-blur conflict: notes.md:20-23 headline change vs absent config key (inactive).
%   - Why the physics losses are constructed but never applied (deliberate ablation vs side effect of
%     the _run_epoch monkey-patch). No note or commit message explains it.
%   - Training hyperparameter conflict: notes.md:25 (2e-4 / 500 / null) vs config.yaml (1.5e-05 /
%     510 / last_model.pt) -- state which produced the reported checkpoint.
%   - Why AMP and TF32 were disabled for exp059.
%   - Why find_unused_parameters=True was enabled in DDP (no comment or note records a reason).
%   - Why 250 MHz, a training anchor, is in eval_off_anchor_mhz with the highest weight (3.0), and
%     which aggregate quantity checkpoints/best_off_anchor_model.pt is selected on.
%   - Capacity-phase mixing probabilities: config.yaml:47-48 (0.65 / 0.4) vs notes.md:12-14 and
%     cursor-handoff.md:33-37 (0 / 0).
%   - Training hardware / software environment and wall-clock time (no committed artifact).
%   - exp059 model parameter count (printed at train_core.py:1424, not stored).
%   - Decision on subsec:aux_guidance_network -- rename to the exp038 impedance surrogate, or
%     implement an exp059-native component. No exp059 component of that name exists.
%   - Full expression of the Stage-2 exceed maximum-excursion term and gap term (optimize.py:313-329).
%   - Provenance of configs/target_impedance.npy (0.8 ohm floor, 49.99999999999999 ohm ceiling).
%   - Stage-2 normalization path conflict: optimize.py:105 (missing directory) vs config.yaml:134;
%     plus the work needed to wire Stage 2 to the exp059 checkpoint.
%   - Reconciliation of the three budget ranges: Stage 2 K=1..25, AL K=3..10, dataset K<=30.
%   - Rationale for the AL MHz strata quotas (250 MHz -> 140, 155 MHz -> 100 per K).
%   - AL budget conflict: docs/active-learning.md sec.3-4 (exp057-era) vs exp059_gp_error.json.
%   - GP score conflict: inference_pool.py:177-179 docstring (UCB with sigma) vs score_mode 'mu'.
%   - Reason the cached GP fit is cleared at the start of every per-K acquire (per_k_acquire.py:66-69
%     shows the clearing; no comment states why).
%   - Reason for the 128.0 AL overlay sample weight and how the value was chosen.
%   - Author's rationale for residual-GP over MC-dropout uncertainty.
%   - AL fine-tune mixing conflict: config_al_finetune.yaml (0.9 / 0.7) vs exp059_gp_error.json
%     (0.65 / 0.4, wins at runtime) vs cursor-handoff.md sec.2 and CLAUDE.md (0.9 / 0.7).
%
% Citation verification queue:
%   - STILL note=VERIFY in refs.bib (bibliographic fields not confirmed): settles2009active,
%     maddison2016concrete, higgins2017beta, shi2019mmvae. These four must be checked before
%     submission.
%   - hensman2013bigdata, titsias2009variational, gardner2018gpytorch, kennedy2001calibration were
%     added to refs.bib earlier; the chapter audit reports full author/year/title fields and no
%     VERIFY note on them.
%   - All 21 cite keys used here resolve to entries in Thesis_report/refs.bib; none is invented.
%     Full list: kingma2013auto, sohn2015conditional, higgins2017beta, wu2018mvae,
%     kutuzova2021defense, shi2019mmvae, perez2018film, kipf2016graph, gilmer2017message,
%     bengio2013straight, jang2016gumbel, maddison2016concrete, gomezbombarelli2018automatic,
%     srinivas2010gp, gal2016dropout, settles2009active, batchal2022batch, gardner2018gpytorch,
%     titsias2009variational, hensman2013bigdata, kennedy2001calibration.
%   - The residual-GP construction is stated in the text as an ASSEMBLY of three literatures
%     (SVGP, GP software, model discrepancy). Do not let editing collapse it into one precedent.
%
% Evidence warnings:
%   - INACTIVE LOSSES: the physics terms (ri/cs/ar) and the spectral anti-blur loss are configured
%     and implemented but never enter the exp059 total loss. The chapter states both explicitly.
%     The model must not be described as physics-informed in its training objective.
%   - subsec:aux_guidance_network has no implementation under that name anywhere in the repository.
%     The label is kept as approved; the subsection says so and describes the nearest artifact
%     (exp038 surrogate_impedance.py, legacy forks only).
%   - Stage 2 (pipelines/latent/optimize.py) defaults to EXPERIMENT = 'exp038_true_multi' and is NOT
%     wired to exp059. No end-to-end inverse-design result may be attributed to the exp059 surrogate.
%   - The STE is applied only on the surrogate branch (surrogate(_ste_topk(occ, K))); the fallback
%     VAE-head branch does not exercise it.
%   - acq_ab_gp_vs_mc is UNCERTAIN and the exp059 acquisition A/B is still an open task. No claim of
%     GP superiority over MC or random appears in this chapter, and none may be added.
%   - use_multiscale_film, the headline architectural change of exp059, is set by a code default in
%     exp059_common.py:43, not by config.yaml.
%   - Stale documents: docs/normalization-and-losses.md (2026-06-22, exp050 Tier A, six anchors,
%     heatmap weight 4.0), docs/training-procedure.md (exp054-exp057 lineage, self-declared caveat
%     that epoch counts/weights/off-anchor MHz are config-specific), docs/active-learning.md
%     (exp057-era budgets), docs/dataset.md (16-anchor claim), root README.md (exp057 as current).
%     None was used as a source for numbers; the body now states this once, in general terms, in
%     Section sec:experimental_setup. Specifics are in the Implementation Index.
%   - REGISTER PASS: invocation mechanics removed from the body (module launch command, torchrun
%     command, the JSON-parsed .yaml config format, the CONFIGURATION-block script convention, the
%     absolute Windows ERF path, run.py/pipeline.py entry-point mechanics). All are preserved in the
%     Implementation Index. Restore them to an appendix, not to the narrative.
%   - Training hardware was deliberately left as \fillin rather than taken from a live shell.
%   - EVIDENCE-BLOCK ERROR: the extracted facts state "the occupancy and impedance decoders see only
%     the shared part". That is wrong for the impedance decoder. vae_multi_input_simple.py:294-296
%     gives dec_in_occ = shared_latent_dim + cond_dim but dec_in_imp = layout_latent_dim + cond_dim
%     + 64, and :174 defines layout_latent_dim = shared_latent_dim + imp_peak_dim. The chapter
%     follows the code. Do not restore the facts-block wording.
%   - The shared=80 statement is at notes.md:15. PhysicsLoss is instantiated at train_core.py:1393
%     and the physics guard is train_core.py:609-613.
%   - Off-anchor evaluation is a misnomer for one of its three frequencies: 250.0 MHz is a training
%     anchor and carries the highest weight (3.0). Only 155.0 and 265.0 are held out.
%   - RATIONALE CLAUSES REMOVED: five "so that / which is required because / is what makes" clauses
%     -- per-anchor normalization, modality dropout, find_unused_parameters, the 128.0 overlay
%     weight, and the per-cycle GP cache clearing -- asserted a mechanism with no ablation, comment
%     or note behind it. They are now either design intent or \fillin. Do not reinstate them as fact.
% ---------------------------------------------------------------
